\documentclass[11pt,a4paper]{amsart}
\usepackage[T1]{fontenc}
\usepackage{amssymb}
\usepackage[margin=1in]{geometry}
\usepackage[hidelinks]{hyperref}
\DeclareFontShape{OT1}{cmr}{bx}{n}{<-> cmbx10}{}
\DeclareMathSizes{8}{8}{8}{8}
\DeclareMathSizes{9}{9}{8}{8}
\DeclareMathSizes{10.95}{10.95}{8}{8}
\urlstyle{same}
\hypersetup{pdftitle={Smoothness of Nash CI surfaces for two disjoint edges}}
\newtheorem{theorem}{Theorem}
\newcommand{\PP}{\mathbb P}
\newcommand{\CC}{\mathbb C}
\title{Smoothness of Nash CI surfaces for two disjoint edges}

\date{}
\begin{document}
\begin{abstract}
For a binary game whose conditional-independence graph consists of two disjoint
edges and any number of isolates, we prove that the Nash conditional-independence
variety is a smooth complex projective surface for generic payoff tables. The
reduced payoff equations are known. We prove full Jacobian rank on their fixed
base loci by stratifying the edge factors and bounding the codimension of
dependence among the normal differentials.
\end{abstract}
\maketitle

\section{Statement and prior work}

Consider $n=r+4$ players with two actions each. The graph specifies the
factorization of the joint strategy. Each player's payoff array is independently
arbitrary on the $2^n$ pure action profiles, and we work over $\CC$.
In tensor coordinates $p_{\mathbf a}$, the Nash conditional-independence
variety is the projective closure of the equal-conditional-payoff rank equations
on the graphical model off the hyperplanes $p_{\mathbf a}=0$ and
$\sum_{\mathbf a}p_{\mathbf a}=0$ \cite{game,totally}. We take the reduced closure.

\begin{theorem}\label{thm:main}
For every integer $r\geq0$, there is a nonempty Zariski open subset of the
affine space of payoff arrays on which the Nash CI variety for two disjoint
edges and $r$ isolated vertices is a smooth complex projective surface.
\end{theorem}

This is the binary surface case of Sendra-Arranz's Question 3.6.1
\cite[p.~196]{thesis}, as isolated in catalog Q220 \cite{catalog}.
Portakal and Sendra-Arranz established smoothness of the one-edge curves
\cite[Theorem~17]{curve}. Their surface work gives the binary defining
equations, the dimension, connectedness and the general-type conclusion
conditional on smoothness \cite{game}. Bouyer, Portakal and Sendra-Arranz
established generic irreducibility of the cluster-graph family
\cite[Theorem~4.4]{totally}. Their smoothness result excludes the tensor-coordinate
and total-mass hyperplanes \cite[Theorem~3.7]{totally}. Theorem~\ref{thm:main}
includes all points of the complex projective closure.

\section{Payoff equations}

The closed graphical model is the Segre embedding of
\[
M=\PP^3_A\times\PP^3_B\times(\PP^1)^r.
\]
Here $A$ and $B$ are $2$ by $2$ matrices of homogeneous coordinates, one for
each edge. We call the two players in a factor its row and column players.
Write $R_0,R_1$ for the row sums in $A$. For each $i$ let $L_i$ be the
linear combination, with coefficients from the row player's payoff table for
action $i$, of entries of row $i$ multiplied by one coordinate from each
other factor. The reduced row-player equation is
\[
f=R_1 L_0-R_0 L_1=0.
\]
The column-player equation is its transpose with a separate payoff array.
The equations on $B$ are analogous. An isolated player's reduced equation is
the difference of the two payoff forms on the other factors and is
multilinear there. These equations follow by cancelling the common
factors in the conditional payoff rank conditions on the indicated chart.
They are also the binary case of the equations in \cite[Proposition~4.3]{totally}.

At any fixed point of the other factors there is a nonzero coordinate
monomial, so the coefficients of $L_0,L_1$ as supported linear forms in
the edge matrix range independently over all such coefficients. For an
isolated player the same observation makes evaluation of its equation
surjective at every projective point. The $n$ equations use independent
payoff blocks. Let $T$ be the full affine parameter space, $Y_t$ the
zero scheme on $M$ at $t\in T$, and $D$ the pullback of the union of the
coordinate and total-mass hyperplanes. Each defining linear form is
nonzero on $M$, so $\dim D\leq \dim M-1$. Off $D$ the original rank
equations agree with the reduced equations.

\section{Base strata and their differentials}

For a single edge factor, the row player's base locus consists of the
two lines where one row is zero and the line where both row sums are
zero. Indeed, evaluation vanishes for every payoff array precisely
when $R_1$ times row $0$ and $R_0$ times row $1$ both vanish.
The column player's lines are obtained by transposition.
This specializes the linear series analyzed in \cite[Lemmas~7--8]{game}.
We partition the factor into the complement of the six lines, their
open parts with intersections removed, and the intersection points.

A same-player intersection occurs when one row is zero and the other
row sums to zero, or the transpose. There are four such points. A
row-column intersection is one of the four single-entry coordinate
points or the checkerboard point represented by
$\left(\begin{smallmatrix}1&-1\\-1&1\end{smallmatrix}\right)$.
A zero row together with both column sums zero forces every entry to
vanish, and the transposed statement is the same. This exhausts the
intersections and rules out triple intersections. The codimension
$c$ and number $b$ of base equations in the factor are
\[
\begin{array}{c|cc}
\text{stratum} & c & b\\ \hline
\text{complement} & 0&0\\
\text{open line} & 2&1\\
\text{same-player point} & 3&1\\
\text{row-column point} & 3&2.
\end{array}
\]
We use products of these strata with all singleton factors. Both $c$
and $b$ add under products.

A base equation vanishes identically as the other factors vary with
its own base point fixed, so its differential has no contribution
from the other factors. It annihilates the tangent of the stratum.
The following images are calculated by varying the base player's payoff block.

\begin{itemize}
\item On an open zero-row line, say row $0$ is zero, $R_1\ne 0$ and
$df=R_1\,dL_0-L_1\,dR_0$, giving the whole two-dimensional conormal.
On the open row-sum line, $df=L_0\,dR_1-L_1\,dR_0$; the two nonzero
rows let us vary $L_0,L_1$ independently, giving the whole conormal.
\item At a same-player point with row $0$ zero and $R_1=0$,
$df=-L_1\,dR_0$ has one-dimensional image since row $1$ is nonzero.
\item At a row-column coordinate point write
$A=\left(\begin{smallmatrix}1&x\\y&z\end{smallmatrix}\right)$
after permuting rows and columns. At $(x,y,z)=(0,0,0)$ the row and
column differential images are $\langle dy,dz\rangle$ and
$\langle dx,dz\rangle$. At the checkerboard point they are
$\langle dx,dy+dz\rangle$ and $\langle dy,dx+dz\rangle$.
In each case the planes are distinct and the two payoff blocks attain
them independently.
\end{itemize}

\section{Generic rank across the strata}

Fix a stratum $S$ with codimension $c$ and $b$ base equations. The
incidence of points of $S$ and payoffs satisfying the $n-b$ nonbase
equations is a vector bundle over $S$. The evaluation from the
trivial payoff bundle to the sum of the $n-b$ equation line bundles
is surjective at each point by the independent blocks and the
definition of $S$. Its kernel is the incidence bundle. Hence the
incidence is smooth of dimension
\[
\dim S+\dim T-(n-b)=\dim T+2-c+b.
\]
Generic smoothness in characteristic zero, applied over the target
$T$, gives a nonempty open set of $t$ where the fibers of the
incidence projection are empty or smooth of dimension $2-c+b$.
For a nondominant stratum we remove the closure of the image.
On a dominant stratum the projection differential is onto over
this open set. The incidence tangent vectors $(v,w)$ satisfy
$F_pv+F_tw=0$ for the nonbase equation derivatives, and $F_t$ is onto.
Thus surjectivity of the projection implies $F_p$ is onto. The
nonbase differentials restricted to the tangent of $S$ are
independent everywhere on the generic fiber.

It remains to check the normal differentials of the base equations.
For a fixed point of the stratum, none of the incidence equations
constrains the base players' payoff blocks. A base differential on
an open line vanishes in codimension two in its block; at a
same-player intersection it vanishes in codimension one. At a
row-column point the bad ordered pairs in the product of the two
distinct two-planes consist of pairs with a zero vector and pairs
both in the common line. Each of these three pieces has codimension
two. The independently surjective differential maps preserve their
codimensions under inverse image.

If a factor contributes an open line or a row-column point, then
$c-b\geq 1$ for the product stratum. The incidence dimension is
at most $\dim T+1$, so the corresponding rank-loss locus, of
relative codimension two, has dimension at most $\dim T-1$.
If a factor contributes a same-player point, then $c-b\geq2$ and
the incidence dimension is at most $\dim T$; the corresponding
bad locus has relative codimension one. Again its dimension is at
most $\dim T-1$. There are finitely many strata and each bad locus
is algebraic on its stratum. The closures of their projections
therefore remove only a proper closed subset of $T$.

Outside the excluded subsets the base differentials are independent.
Different edge factors occupy different cotangent summands, and
within each factor we have excluded its dependence locus. They
annihilate the stratum tangent, while the nonbase differentials
restrict to an independent set there. All $n$ differentials are
independent at each point of $Y_t$. The Jacobian criterion gives
smoothness and pure dimension $\dim M-n=2$ whenever $Y_t$ is nonempty.

\section{Nonemptiness and closure}

Write $h_A,h_B,z_1,\ldots,z_r$ for the hyperplane classes and
$Z=\sum_i z_i$. The first four equation classes are two copies each
of $2h_A+h_B+Z$ and $h_A+2h_B+Z$. The isolated-player classes
are $h_A+h_B+Z-z_i$. Their product in the Chow ring is nonzero.
For $r=0$ it contains $h_A^2h_B^2$, and for $r=1$ it contains
$h_A^3h_B^2$. For $r\geq2$ choose the other singleton class in
each factor by a cyclic shift to get
$h_A^2h_B^2\prod_i z_i$. These monomials survive the truncations
of the Chow ring and every coefficient is nonnegative. A
nowhere-vanishing section of the direct sum of the equation
line bundles would give zero top Chern class. Thus $Y_t$ is
nonempty. Nonemptiness also follows from \cite[Theorem~5.4]{totally}.

A stratum with a base equation has $c-b\geq1$ and so its generic
fiber has dimension at most one. On the zero-base stratum, restrict
the incidence bundle to $D$. The dimension of the resulting
incidence is at most $\dim T+1$, so its generic fiber has dimension
at most one by the fiber dimension theorem. Thus no component
of $Y_t$ lies in $D$. Since $Y_t$ is reduced and the Segre
embedding is closed, the defining projective closure is $Y_t$ itself.
This proves Theorem~\ref{thm:main}.

\section*{Scope}
The result concerns the projective complex variety. The positive
real equilibrium locus and its semialgebraic boundary are separate
questions. The proof controls this binary two-edge family. For
irreducibility we refer to \cite[Theorem~4.4]{totally}.

\section*{Review and disclosure}
Internally reviewed, not peer reviewed. Under human direction,
ChatGPT Astra 6.2 (OpenAI) generated the discovery proof and its
research packet. Claude Opus 5.5 (Anthropic) assisted with earlier
mathematical and literature audits. OpenAI Codex performed the mathematical review and
revised the paper. GPT-6.2 Sol assisted with literature searches and
release preparation. DannyExperiments maintains the repository and
is the preferred citation identity. No AI system is proposed as an
author.

\begin{thebibliography}{9}
\bibitem{thesis} J. Sendra-Arranz, \emph{Algebraic Methods in Reconstruction of
Varieties and Game Theory}, doctoral dissertation, Universit\"at T\"ubingen (2025), \S 3.6.
\href{https://publikationen.uni-tuebingen.de/xmlui/handle/10900/160061}{Institutional record}.
\bibitem{catalog} Open Problems Catalog, Q220.
\href{https://github.com/10m-u/open-problems-catalog/blob/ca8d4d44ba3db527f19f28928c240168f545ecfe/problems/game_theory/noncooperative-games-equilibrium-theory.md\#q220}{Pinned entry}.
\bibitem{curve} I. Portakal and J. Sendra-Arranz, Nash conditional
independence curve, \emph{J. Symbolic Comput.} 122 (2024), 102255.
\href{https://doi.org/10.1016/j.jsc.2023.102255}{Journal article}.
\bibitem{game} I. Portakal and J. Sendra-Arranz, Game theory of undirected
graphical models, \emph{J. Algebra} 666 (2025), 574--606.
\href{https://doi.org/10.1016/j.jalgebra.2024.12.004}{Journal article}.
\bibitem{totally} M. Bouyer, I. Portakal and J. Sendra-Arranz, Totally
mixed conditional independence equilibria of generic games (2025).
\href{https://arxiv.org/abs/2511.11467v1}{arXiv:2511.11467v1}.
\end{thebibliography}
\end{document}
